%% ============================================================================
%% Fichier  : annexes/Ann_E_bibliographie.tex
%% Titre    : Bibliographie commentée
%% Version  : 1.0 -- novembre 2026
%% ============================================================================
%% Alignée sur les annexes des autres cours de l'auteur (ACSI, INF3226).
%% ============================================================================

\chapter{Bibliographie Commentée}
\label{ann:biblio}

Pour chaque partie, quelques lectures, avec ce qu'elles apportent. Les références
complètes figurent dans la bibliographie du livre. Les textes de l'auteur sont des
prépublications déposées sur l'IACR ePrint~; aucun travail non publié n'est cité.

\section*{Pour commencer}
\begin{itemize}
    \item Casey~\cite{casey2004}~: le manuel de référence sur la preuve numérique et la criminalité informatique.
    \item Pollitt~\cite{pollitt2010history} et Garfinkel~\cite{garfinkel2010next}~: l'histoire de la discipline et ses défis.
\end{itemize}

\section*{Partie I · Fondements}
\begin{itemize}
    \item Kuhn~\cite{kuhn1962}, Popper~\cite{popper1973}, Foucault~\cite{foucault1969}~: paradigmes, réfutation, régimes de vérité.
    \item Locard~\cite{locard1920}~: l'échange de traces, transposé au numérique.
    \item Shannon~\cite{shannon1948}, Dolev et Yao~\cite{dolevyao1983}~: l'information et l'adversaire.
    \item Abbott~\cite{abbott1988}~: la formation d'une profession.
    \item Le Trilemme CRO~\cite{minka2025cro}~: le fil rouge du cours.
\end{itemize}

\section*{Partie II · Processus et méthodologies}
\begin{itemize}
    \item DFRWS~\cite{dfrws2001}~: le premier modèle de processus.
    \item ISO/IEC~27037, 27041, 27042, 27043~\cite{iso27037,iso27041_2015,iso27042_2015,iso27043_2015}~; NIST SP~800-86 et 800-101~\cite{nist80086,nist800101r1}~; RFC~3227~\cite{rfc3227}~; guide de l'ACPO~\cite{acpo2012}.
    \item Heuer~\cite{heuer1999}~: l'analyse des hypothèses concurrentes.
    \item Carrier~\cite{carrier2005}~: les systèmes de fichiers, pour comprendre ce que l'on acquiert.
\end{itemize}

\section*{Partie III · Cadre normatif et déontologie}
\begin{itemize}
    \item Convention de Budapest~\cite{budapest2001}, Convention de Malabo~\cite{malabo2014}, RGPD~\cite{rgpd2016}, CLOUD Act~\cite{cloudact2018}.
    \item Lois n°~2010/012, 2010/013 et 2024/017~\cite{cameroun2010a,cameroun2010b,cameroun2024donnees}.
    \item Dror~\cite{dror2020cognitive}~: les biais cognitifs de l'expert~; Milne et Bull~\cite{milne1999investigative}~: l'entretien d'enquête~; Becker~\cite{becker1963}~: l'étiquetage.
\end{itemize}

\section*{La suite}
L'UE SEC4052 prolonge ces lectures vers la forensique opérationnelle et les fondements cryptographiques de la preuve.
